% ============================================================
%  RF Toolbox Textbook — shared preamble
%  Compiles with pdflatex (utf8). Figures are PDFs from rsvg-convert.
% ============================================================
\usepackage[utf8]{inputenc}
\usepackage[T1]{fontenc}
\usepackage{lmodern}
\usepackage[a4paper,margin=1in]{geometry}
\usepackage{amsmath,amssymb}
\usepackage{graphicx}
\usepackage[export]{adjustbox}
\usepackage{booktabs}
\usepackage{array}
\usepackage{longtable}
\usepackage{enumitem}
\usepackage{xcolor}
\usepackage[most]{tcolorbox}
\usepackage{titlesec}
\usepackage{sectsty}
\usepackage{newunicodechar}
\usepackage{textcomp}
\usepackage[hidelinks]{hyperref}
\hypersetup{colorlinks=true,linkcolor=rfpurple,urlcolor=rfaccent,citecolor=rfpurple,
  pdftitle={RF and Microwave Engineering — An Interactive Textbook},
  pdfauthor={RF Toolbox (rftoolbox.ca)}}

% ---- theme colours (match the site) ----
\definecolor{rfpurple}{HTML}{5533AA}
\definecolor{rfaccent}{HTML}{AA77FF}
\definecolor{rffill}{HTML}{F0EEFF}
\definecolor{rffill2}{HTML}{F8F6FF}
\definecolor{rfamber}{HTML}{CC8800}
\definecolor{rfamberfill}{HTML}{FFFBEE}

% ---- part / chapter / section styling ----
\allsectionsfont{\sffamily}
\chapterfont{\sffamily\color{rfpurple}}
\titleformat{\section}{\sffamily\large\bfseries\color{rfpurple}}{\thesection}{0.7em}{}
\titleformat{\subsection}{\sffamily\bfseries\color{rfpurple!85}}{\thesubsection}{0.6em}{}

% ============================================================
%  Unicode -> LaTeX mappings (covers every non-ASCII char in the corpus)
% ============================================================
\newunicodechar{—}{\textemdash}
\newunicodechar{–}{\textendash}
\newunicodechar{·}{\ensuremath{\cdot}}
\newunicodechar{−}{\ensuremath{-}}
\newunicodechar{×}{\ensuremath{\times}}
\newunicodechar{÷}{\ensuremath{\div}}
\newunicodechar{±}{\ensuremath{\pm}}
\newunicodechar{≈}{\ensuremath{\approx}}
\newunicodechar{≤}{\ensuremath{\leq}}
\newunicodechar{≥}{\ensuremath{\geq}}
\newunicodechar{≪}{\ensuremath{\ll}}
\newunicodechar{≠}{\ensuremath{\neq}}
\newunicodechar{∝}{\ensuremath{\propto}}
\newunicodechar{∞}{\ensuremath{\infty}}
\newunicodechar{√}{\ensuremath{\surd}}
\newunicodechar{→}{\ensuremath{\rightarrow}}
\newunicodechar{∑}{\ensuremath{\textstyle\sum}}
\newunicodechar{°}{\ensuremath{{}^{\circ}}}
\newunicodechar{µ}{\ensuremath{\mu}}
\newunicodechar{ℓ}{\ensuremath{\ell}}
\newunicodechar{è}{\`e}
\newunicodechar{α}{\ensuremath{\alpha}}
\newunicodechar{β}{\ensuremath{\beta}}
\newunicodechar{γ}{\ensuremath{\gamma}}
\newunicodechar{δ}{\ensuremath{\delta}}
\newunicodechar{ε}{\ensuremath{\varepsilon}}
\newunicodechar{η}{\ensuremath{\eta}}
\newunicodechar{θ}{\ensuremath{\theta}}
\newunicodechar{λ}{\ensuremath{\lambda}}
\newunicodechar{μ}{\ensuremath{\mu}}
\newunicodechar{π}{\ensuremath{\pi}}
\newunicodechar{σ}{\ensuremath{\sigma}}
\newunicodechar{τ}{\ensuremath{\tau}}
\newunicodechar{φ}{\ensuremath{\varphi}}
\newunicodechar{χ}{\ensuremath{\chi}}
\newunicodechar{ω}{\ensuremath{\omega}}
\newunicodechar{Γ}{\ensuremath{\Gamma}}
\newunicodechar{Δ}{\ensuremath{\Delta}}
\newunicodechar{Π}{\ensuremath{\Pi}}
\newunicodechar{Σ}{\ensuremath{\Sigma}}
\newunicodechar{Ω}{\ensuremath{\Omega}}
\newunicodechar{¹}{\textsuperscript{1}}
\newunicodechar{²}{\textsuperscript{2}}
\newunicodechar{³}{\textsuperscript{3}}
\newunicodechar{⁴}{\textsuperscript{4}}
\newunicodechar{⁵}{\textsuperscript{5}}
\newunicodechar{⁺}{\textsuperscript{+}}
\newunicodechar{⁻}{\textsuperscript{$-$}}
\newunicodechar{₀}{\textsubscript{0}}
\newunicodechar{₁}{\textsubscript{1}}
\newunicodechar{₂}{\textsubscript{2}}
\newunicodechar{₃}{\textsubscript{3}}
\newunicodechar{₄}{\textsubscript{4}}
\newunicodechar{₅}{\textsubscript{5}}
\newunicodechar{₊}{\textsubscript{+}}
\newunicodechar{₋}{\textsubscript{$-$}}
\newunicodechar{ₚ}{\textsubscript{p}}
\newunicodechar{ᵣ}{\textsubscript{r}}

% ============================================================
%  Pedagogical environments
% ============================================================
\newtcolorbox{objectives}{colback=rffill,colframe=rfaccent,boxrule=0pt,
  leftrule=3pt,arc=2pt,left=10pt,right=10pt,top=6pt,bottom=6pt,
  title={\sffamily\bfseries\footnotesize\MakeUppercase{Learning Objectives}},
  coltitle=rfpurple,fonttitle=\bfseries}

\newtcolorbox{keytakeaways}{colback=rfamberfill,colframe=rfamber,boxrule=0pt,
  leftrule=3pt,arc=2pt,left=10pt,right=10pt,top=6pt,bottom=6pt,
  title={\sffamily\bfseries\footnotesize\MakeUppercase{Key Takeaways}},
  coltitle=rfamber!80!black,fonttitle=\bfseries}

\newtcolorbox{workedexample}[1][Worked Example]{colback=white,colframe=rfpurple!55,
  boxrule=0.6pt,arc=2pt,left=10pt,right=10pt,top=6pt,bottom=6pt,
  title={\sffamily\bfseries\footnotesize\MakeUppercase{#1}},coltitle=white,
  colbacktitle=rfpurple,fonttitle=\bfseries}

\newtcolorbox{conceptcheck}{colback=rffill2,colframe=rfaccent,boxrule=0.6pt,arc=2pt,
  left=10pt,right=10pt,top=6pt,bottom=6pt,
  title={\sffamily\bfseries\footnotesize\MakeUppercase{Concept Check}},
  coltitle=rfpurple,fonttitle=\bfseries}

\newenvironment{exercises}{\subsection*{Exercises}\begin{enumerate}[leftmargin=1.6em,itemsep=4pt]}{\end{enumerate}}
\newtcolorbox{answer}{colback=rffill2,colframe=rfaccent,boxrule=0pt,leftrule=2.5pt,
  arc=1pt,left=8pt,right=8pt,top=3pt,bottom=3pt,before skip=3pt,after skip=6pt,
  fontupper=\small}

% figure helper: keep figures within the text block, cap height
\newcommand{\rffig}[2]{%
  \begin{figure}[htbp]\centering
  \includegraphics[max width=\linewidth,max totalheight=0.42\textheight]{#1}%
  \ifx\relax#2\relax\else\caption{#2}\fi
  \end{figure}}
